최근 대규모 언어 모델(LLM)을 활용한 멀티 에이전트 시스템이 복잡한 문서 생성 작업에 도입되고 있으나, 에이전트 간의 협업을 제어하는 오케스트레이션 구조에 따른 성능 차이에 대한 정량적 분석은 부족한 실정이다. 본 연구는 오케스트레이션 방식이 생성 문서의 품질과 자원 효율성에 미치는 영향을 분석하여, 작업 성격에 최적화된 제어 구조를 제안하는 것을 목적으로 한다. 이를 위해 고정된 그래프 기반의 Code 방식, 동적 라우팅 기반의 LLM 방식, 그리고 계층적 제어를 수행하는 Hybrid 방식의 세 가지 오케스트레이션 구조를 정의하고 구현하였다. 산업 현장의 정형 보고서 작성 및 창의적 글쓰기 작업을 대상으로 토큰 소비량, 실행 시간과 같은 효율성 지표와 LLM-as-Judge 기반의 품질 점수를 측정하여 비교 분석을 수행하였다. 실험 결과, 각 구조는 자원 효율성과 생성 품질 간의 뚜렷한 트레이드오프 관계를 보였으며, 제어 메커니즘의 유연성에 따라 재시도 횟수와 최종 품질 점수가 상이하게 나타남을 확인하였다. 본 연구는 오케스트레이션 구조별 제어 메커니즘을 정립하고 정량적 성능 지표를 제시함으로써, 작업의 특성에 따른 최적의 구조 선택 기준을 제공한다는 점에서 학술적 및 실무적 기여를 갖는다.